\begin{enumerate}

\item La tensión circunferencial es $pR/t$.

\item La longitudinal es la mitad.

\item Con $\sigma_1 = \sigma_\theta$ y $\sigma_2 = \sigma_x = \sigma_\theta/2$, $\sigma_v = \sigma_\theta\sqrt{3}/2$.

\end{enumerate}

\noindent\textbf{Operaciones y resultados intermedios.}

\begin{enumerate}

\item Tensión circunferencial.

\[\fitmath{y_{1}=\frac{0.06\cdot 1900}{1.6}\approx 71{,}25\,\mathrm{MPa}}\]

\item Tensión longitudinal.

\[\fitmath{y_{2}=\frac{0.06\cdot 1900}{2\cdot 1{,}6}\approx 35{,}625\,\mathrm{MPa}}\]

\item Tensión de von Mises.

\[\fitmath{y_{3}=\frac{\frac{\sqrt{3}}{2}\cdot 0{,}06\cdot 1900}{1.6}\approx 61{,}70431\,\mathrm{MPa}}\]

\end{enumerate}
